% rotrta.tex — IEEE conference style
% Compile with: pdflatex rotrta.tex (twice for refs)
\documentclass[conference]{IEEEtran}

\usepackage{cite}
\usepackage{amsmath,amssymb}
\usepackage{graphicx}
\usepackage{xcolor}
\usepackage{booktabs}
\usepackage{listings}
\usepackage{url}
\usepackage{microtype}
\usepackage{siunitx}
\usepackage{multirow}
\usepackage{array}

% Code listing style
\lstset{
  basicstyle=\ttfamily\scriptsize,
  keywordstyle=\bfseries\color{blue!70!black},
  commentstyle=\itshape\color{green!50!black},
  stringstyle=\color{red!60!black},
  breaklines=true,
  frame=single,
  numbers=left,
  numberstyle=\tiny,
  xleftmargin=1.5em,
  framexleftmargin=1.5em,
}

\hyphenation{Cat-Whisper Vul-kan FAISS}

% ============================================================================
\begin{document}

\title{ROTRTA: Ridiculous Optimizations Through Ridiculous Test Assertions ---
       A Methodology for AI-Driven Performance Engineering, with Application
       to GPU Vector Similarity Search}

\author{
  \IEEEauthorblockN{IreGaddr}
  \IEEEauthorblockA{CatWhisper Project\\
    \texttt{iregaddr@gmail.com}}
}

\maketitle

% ============================================================================
\begin{abstract}
We introduce \emph{Ridiculous Optimizations Through Ridiculous Test Assertions}
(ROTRTA), a test-driven performance engineering methodology in which
latency assertions are deliberately set beyond a system's current capability,
anchored to the measured performance of a reference state-of-the-art system.
Tests are written first and are contractually immutable: they may never be
weakened to accommodate a lagging implementation.  This asymmetry forces
genuine optimization rather than satisficing, and provides AI-agent optimizers
with a binary, reproducible oracle.  We demonstrate ROTRTA on CatWhisper, a
cross-vendor GPU vector similarity query library built on Vulkan.  Starting
from a manually optimized baseline that remained \SI{2.4}{\times} slower than
FAISS-GPU across three benchmark configurations, a single ROTRTA session
driven by an autonomous AI agent produced a stack of six coordinated
optimizations --- Structure-of-Arrays database layout, doubled warp-tile size,
Vulkan timeline semaphores, userspace spin-poll synchronization, AVX-512
fp16 conversion, and zero-copy mapped I/O --- that together reduced
single-query latency by up to \SI{18.9}{\times} and caused CatWhisper to
\emph{beat FAISS-GPU on all three benchmark dimensions} simultaneously.
We argue that ROTRTA is a broadly applicable discipline that makes latent
performance budget visible, eliminates premature satisficing, and is
uniquely well-suited to AI-agent-driven development because it reduces an
open-ended optimization problem to a well-posed constraint satisfaction
problem.
\end{abstract}

\begin{IEEEkeywords}
test-driven development, performance engineering, vector similarity query,
GPU computing, Vulkan, FAISS, AI agents, benchmark-driven optimization
\end{IEEEkeywords}

% ============================================================================
\section{Introduction}
\label{sec:intro}

High-performance systems development has long grappled with a fundamental
tension: performance requirements are continuous and comparative, yet
traditional test frameworks treat correctness as binary.  A function either
returns the right answer or it does not; but whether it returns it \emph{fast
enough} depends on hardware, workload, and what competitors have achieved.
Standard benchmarking practices measure after the fact.  Profiling-driven
optimization follows an iterative, ad hoc loop --- profile, identify hot spots,
optimize, repeat.  This loop is well understood for human developers.  It is
poorly matched to the working style of autonomous AI agents.

AI coding agents, in their current form, excel at implementing well-specified
transformations and exploring large, locally constrained search spaces.  They
are less effective at formulating their own performance targets, recognizing
satisficing behavior, or maintaining long-horizon optimization intent across
many discrete decisions.  Given a vague directive such as ``make the search
function faster,'' an agent may apply superficial improvements --- reordering
memory accesses, inlining a function --- and declare success.  Given a concrete,
measurable, binary criterion --- ``the median single-query latency for
10{,}000 vectors at dimension 128 must be strictly less than \SI{0.065}{\ms}''
--- the agent has an unambiguous oracle and cannot satisfice prematurely.

ROTRTA (\textbf{R}idiculous \textbf{O}ptimizations \textbf{T}hrough
\textbf{R}idiculous \textbf{T}est \textbf{A}ssertions) formalizes this
observation into a methodology:

\begin{enumerate}
  \item \textbf{Measure the reference.}  Establish the performance of the
        state-of-the-art system on the target hardware and workload.
  \item \textbf{Write immutable assertions.}  Encode performance targets as
        standard unit-test assertions, set at the reference or slightly
        below (25--50\% reduction from the current implementation).  The
        tests are written before the corresponding optimizations exist.
  \item \textbf{Lock the tests.}  Once written, test assertions are
        contractually immutable.  They may never be weakened or deleted
        to make a failing implementation appear correct.
  \item \textbf{Implement until the tests pass.}  Any optimization
        technique that causes the assertions to pass is valid.
        Optimizations that address symptoms without curing root causes
        will reveal themselves when other assertions remain red.
\end{enumerate}

In this paper we present ROTRTA as a discipline (Section~\ref{sec:rotrta}),
describe CatWhisper --- a Vulkan-based GPU vector search library and the
system against which ROTRTA was applied --- contrasting its design with
FAISS (Section~\ref{sec:catwhisper}), trace the full optimization arc from
baseline to ROTRTA-passing implementation (Section~\ref{sec:application}),
and analyze the experimental results (Section~\ref{sec:results}).  We then
discuss broader implications for AI-assisted software engineering
(Section~\ref{sec:discussion}) and conclude (Section~\ref{sec:conclusion}).

% ============================================================================
\section{Background}
\label{sec:background}

\subsection{Vector Similarity Search}

Dense vector similarity search is a core primitive in modern information
retrieval, recommendation systems, and retrieval-augmented generation
(RAG) pipelines.  Given a database $\mathcal{D} = \{v_1, \ldots, v_n\}
\subset \mathbb{R}^d$ and a query $q \in \mathbb{R}^d$, the $k$-nearest
neighbor ($k$-NN) problem finds the $k$ elements of $\mathcal{D}$ minimizing
a distance metric $\delta(q, v_i)$.  For an exact flat index, this requires
evaluating all $n$ distances: $O(nd)$ arithmetic operations per query.

\subsection{FAISS}

Facebook AI Similarity Search (FAISS)~\cite{johnson2019billion} is the
de facto reference library for vector similarity search.  It provides both
CPU and GPU backends, the latter accelerated by CUDA.  FAISS implements
multiple index families: IndexFlat (exact brute-force), IndexIVFFlat
(inverted file), IndexIVFPQ (product-quantized IVF), and HNSW
(hierarchical navigable small world graphs).  GPU-accelerated IndexFlat
achieves throughputs on the order of hundreds of millions of distance
computations per second on modern NVIDIA hardware.  As of FAISS 1.13.2
on an RTX~4080 Laptop GPU, single-query L2 search latency at $n=10{,}000$,
$d=128$, $k=10$ is \SI{0.065}{\ms}; at $n=100{,}000$, $d=128$ it is
\SI{0.350}{\ms}; and at $n=100{,}000$, $d=256$ it is \SI{0.601}{\ms}.

\subsection{Vulkan Compute}

Vulkan~\cite{vulkan2024spec} is a low-overhead, cross-vendor GPU API
standardized by the Khronos Group.  Vulkan compute shaders, compiled to
SPIR-V, execute on any Vulkan-capable GPU: NVIDIA, AMD, Intel, ARM Mali,
and Apple Silicon (via MoltenVK).  Vulkan 1.2 introduced timeline
semaphores --- monotonically-increasing 64-bit counters that eliminate the
reset-before-reuse requirement of binary fences and enable userspace polling
without kernel transitions on drivers that implement GPU-mapped semaphore
surfaces.

\subsection{AI Agents in Software Engineering}

Large language model (LLM) agents have been applied to code generation,
debugging, and refactoring~\cite{chen2021evaluating}.  Their utility for
low-level performance optimization, however, has been less studied.
Performance optimization requires hardware intuition, awareness of micro-
architectural bottlenecks, and the ability to compose individually modest
improvements into coordinated wins --- capabilities that benefit substantially
from a concrete, measurable success criterion.

% ============================================================================
\section{The ROTRTA Methodology}
\label{sec:rotrta}

\subsection{Motivation: The Satisficing Problem}

Optimization without a quantitative target converges to local minima.  A
developer (or agent) who achieves a 20\% speedup has no mechanism to know
whether 80\% was achievable.  Profiling can identify hot spots but does not
indicate how much improvement is theoretically available.  Conventional TDD
enforces functional correctness but is silent on performance: a function that
takes one second is as ``correct'' as one that takes one microsecond.

The ROTRTA insight is to close this loop by encoding performance targets as
first-class, immutable test assertions, set to match or exceed a reference
implementation.  The choice of a reference that is marginally better than
current performance (25--50\% reduction) ensures the target is achievable
but requires genuine effort.  The choice of a reference that represents the
state of the art (e.g., FAISS-GPU) ensures the target is meaningful.

\subsection{Formal Definition}

Let $L_\text{curr}$ denote the current median latency of a system under
benchmark $B$ on hardware $H$.  Let $L_\text{ref}$ denote the latency of
the reference state-of-the-art system under the same $B$ on $H$.  A ROTRTA
assertion is:
\begin{equation}
  \text{assert}\bigl(L_\text{impl} < L_\text{ref} - \epsilon\bigr)
  \label{eq:rotrta}
\end{equation}
where $\epsilon \geq 0$ is a margin chosen such that
\begin{equation}
  \frac{L_\text{curr} - (L_\text{ref} - \epsilon)}{L_\text{curr}}
  \;\in\; [0.25,\; 0.50]
  \label{eq:budget}
\end{equation}
i.e., the assertion requires a 25--50\% reduction relative to the current
implementation.  $\epsilon$ is typically set so that $L_\text{ref} - \epsilon$
is \SI{1}{\ms} or less below the measured reference, creating a margin that
accounts for measurement noise while still requiring the implementation to
genuinely match the reference class.

\subsection{The Immutability Invariant}

Once written, a ROTRTA assertion must never be modified to be less strict.
This invariant is enforced by treating the test file as append-only after the
initial commit.  The justification is behavioral: if an agent (or developer)
can weaken an assertion, it will do so whenever optimization becomes difficult.
The immutability invariant forces the question back to the implementation.

\subsection{Hardware Agnosticism}

ROTRTA targets are expressed in absolute time units (milliseconds) rather than
relative speedups, but must be calibrated to the target hardware class.  For
this work, calibration was performed on an NVIDIA RTX~4080 Laptop GPU running
FAISS-GPU 1.13.2.  The resulting thresholds are specific to that hardware
class but remain reproducible within it.

\subsection{Relationship to Test-Driven Development}

ROTRTA shares TDD's test-first discipline and immutability of the test
contract.  It differs in three respects: (1) assertions encode
\emph{performance} rather than functional correctness; (2) the reference
values are derived from external benchmarks rather than specifications;
and (3) the optimization space is continuous and hardware-sensitive rather
than discrete and portable.  ROTRTA does not replace TDD; it complements
it by adding a performance layer to the test suite.

\subsection{Suitability for AI Agents}

ROTRTA is particularly well-suited to AI agents for three reasons.

\textbf{Binary oracle.}  An agent interacting with a build-and-test loop
requires an unambiguous success signal.  ROTRTA provides exactly this: the
test suite either passes or fails.  The agent need not form its own judgment
about whether a speedup is ``sufficient.''

\textbf{Composable improvements.}  A single ROTRTA target may require
multiple coordinated changes.  An agent can decompose the problem into
candidate optimizations, apply them incrementally, and use the test oracle
to validate each step.  The binary signal guides search without requiring
the agent to maintain a mental model of cumulative performance impact.

\textbf{Prevention of premature satisficing.}  The immutability invariant
means an agent cannot exit the optimization loop by rationalizing that
``this is fast enough.''  The loop terminates only when the assertion passes.

% ============================================================================
\section{CatWhisper: System Design and Contrast with FAISS}
\label{sec:catwhisper}

CatWhisper is a C++ GPU vector similarity query library targeting
production use.  Its primary design axis --- cross-vendor Vulkan compute
rather than CUDA --- leads to a set of architectural decisions that differ
materially from FAISS, and which constitute both constraints and opportunities
for optimization.

\subsection{Vendor Neutrality via Vulkan}

FAISS's GPU backend is built on CUDA and NVIDIA cuBLAS, limiting it to
NVIDIA hardware.  CatWhisper uses Vulkan 1.2 compute shaders compiled from
GLSL to SPIR-V, and is tested on NVIDIA, AMD, and (via MoltenVK) Apple
Silicon.  Vulkan's lower-level API surface introduces more programmer
responsibility --- explicit synchronization, explicit memory type selection,
explicit command buffer management --- but eliminates the CUDA driver model's
implicit stream management overhead in latency-sensitive paths.

\subsection{Memory Layout: SoA vs.\ AoS}

FAISS stores database vectors in Array-of-Structures (AoS) order:
\begin{equation}
  \text{database}[\text{vector\_idx} \times d + \text{dim\_idx}]
\end{equation}
Under AoS layout, when $T$ threads in a warp each process a different vector,
thread $t$ reads dimension 0 of vector $t$, located at byte offsets
$0, d \cdot s_\text{fp16}, 2d \cdot s_\text{fp16}, \ldots$  These are strided
accesses: each warp-level read pulls $T$ separate cache lines.

CatWhisper uses Structure-of-Arrays (SoA) order:
\begin{equation}
  \text{database}[\text{dim\_idx} \times C + \text{vector\_idx}]
\end{equation}
where $C \geq n$ is the allocated capacity (stride).  When thread $t$ reads
dimension $d$ of its assigned vector at position $\text{base} + t$, it
accesses $\text{database}[d \times C + \text{base} + t]$.  Consecutive
threads differ by one vector index and thus by two bytes (fp16).  All 32
threads in a warp read a contiguous 64-byte cache line --- one read per
dimension step, regardless of the number of threads.  This is the maximum
possible coalescing efficiency.

The SoA layout requires a transposition pass during \texttt{add()}, which
CatWhisper implements as a dedicated compute shader
(\texttt{transpose\_add.comp}) dispatched on the GPU, adding no CPU cost.

\subsection{Data Type: fp16 Storage, fp32 Accumulation}

Both CatWhisper and FAISS support fp16 (half-precision) database vectors.
CatWhisper stores vectors as \texttt{float16\_t} in the SoA buffer and
promotes each value to \texttt{float} before accumulation within the shader:
\begin{lstlisting}[language=C,caption={fp16-to-fp32 accumulation in GLSL}]
float diff = float(database[d*cap+pos])
           - float(query[d]);
sum += diff * diff;
\end{lstlisting}
This halves the database memory footprint, improving cache utilization, and
cuts memory bandwidth by 2$\times$ relative to fp32 storage --- the dominant
cost at large $n$.

\subsection{Fused Distance and Top-K Shader}

FAISS separates distance computation and top-k selection into distinct kernel
invocations, requiring an intermediate distance buffer of size $n$ and a
global memory barrier between passes.  CatWhisper fuses both operations into
a single shader (\texttt{distance\_topk\_fused.comp}):

\begin{enumerate}
  \item Each workgroup of 1{,}024 threads processes a tile of $C_t = 2{,}048$
        vectors, computing distances and writing them into a shared-memory
        sort buffer (\texttt{sort\_buf[2048]}).
  \item An ascending bitonic sort over the 2{,}048-element shared buffer
        is performed in-place, with each thread executing two comparisons
        per stage.
  \item The top-$k$ entries (the first $k$ elements of the sorted buffer)
        are written to a result buffer.
\end{enumerate}

This fusion eliminates the intermediate distance buffer, reduces global memory
traffic by $(n - k \cdot G)$ floats per query (where $G$ is the number of
workgroups), and avoids a global barrier.  The CPU performs a final $O(k G
\log k)$ max-heap merge over the $G \cdot k$ partial results.

\subsection{Synchronization: Timeline Semaphores and Userspace Spin-Poll}

FAISS-GPU uses CUDA stream synchronization (\texttt{cudaStreamSynchronize}),
which issues a kernel and blocks the CPU until the GPU signals completion.
On NVIDIA hardware, the CUDA driver maps a semaphore surface into the
process address space, allowing \texttt{cudaStreamSynchronize} to poll in
userspace without a syscall under favorable conditions.

Vulkan's traditional synchronization primitive, \texttt{VkFence}, requires
\texttt{vkResetFences} before each reuse (a device-visible operation with
$\sim$1--2\,\si{\micro\second} latency) and
\texttt{vkWaitForFences} with a kernel-visible blocking wait (5--15\,\si{\micro\second}
wake-up jitter on Linux).

CatWhisper replaces the fence with a Vulkan 1.2 timeline semaphore ---
a monotonically increasing 64-bit counter that signals after each
\texttt{vkQueueSubmit} without requiring a reset --- and polls it with
\texttt{vkGetSemaphoreCounterValue} in a userspace busy-wait loop.  On
NVIDIA's proprietary driver, \texttt{vkGetSemaphoreCounterValue} reads a
GPU-mapped host-visible surface that the GPU updates on completion,
resulting in sub-microsecond detection latency with no kernel transitions.
A spin limit (\texttt{SPIN\_LIMIT = 2{,}000{,}000} iterations) prevents
livelock: if the semaphore has not signaled within the limit, the code falls
back to \texttt{vkWaitSemaphores} (blocking wait).

The combination eliminates two sources of latency present in the fence-based
path: the reset operation and the kernel sleep/wake round-trip.

\subsection{Zero-Copy Query and Result Paths}

\textbf{Query upload.}  A naive implementation allocates a staging buffer,
copies the query vector into it, issues a \texttt{vkCmdCopyBuffer}, and
submits and waits for the transfer --- a full GPU round-trip per query.
CatWhisper allocates the query buffer with
\texttt{VMA\_MEMORY\_USAGE\_GPU\_TO\_CPU} (HostCoherent) and maintains a
persistent mapped pointer.  Writing a query is a single \texttt{memcpy}
from host memory to the mapped buffer, with no Vulkan call, no staging
buffer, and no fence.

\textbf{Result download.}  Distance and index result buffers are allocated
as \texttt{VMA\_MEMORY\_USAGE\_GPU\_TO\_CPU} (HostReadback, host-cached
system RAM).  After the GPU writes results, the CPU reads directly through
the mapped pointer without a \texttt{vkCmdCopyBuffer} or fence wait.
Host-cached system RAM is read efficiently by the CPU's prefetch hardware,
comparable in bandwidth to reading from DRAM.

\textbf{Query vector conversion.}  When fp16 mode is active, the query
vector must be converted from \texttt{float32} (the public API type) to
\texttt{float16} before the GPU shader can consume it.  CatWhisper performs
this conversion via SIMD intrinsics on the CPU during the \texttt{memcpy}
into the mapped query buffer.  The conversion path uses AVX-512 when
available (\texttt{\_mm512\_cvtps\_ph}, 16 elements per cycle), falls back
to F16C (\texttt{\_mm256\_cvtps\_ph}, 8 elements per cycle), and finally
to a scalar implementation.  On an AMD Ryzen~9 7945HX with AVX-512,
converting a 256-element vector requires approximately 16 SIMD instructions.

\subsection{Persistent Reusable Command Buffer}

FAISS launches CUDA kernels directly; each launch queues work for the stream.
Vulkan requires explicit command buffer recording.  A naive implementation
records a new command buffer for every search call: allocate, begin (with
\texttt{ONE\_TIME\_SUBMIT\_BIT}), bind pipeline, bind descriptor sets, push
constants, dispatch, end, submit, wait, free.  This recording overhead is
\SI{5}{\micro\second}--\SI{15}{\micro\second} on a modern system.

CatWhisper records the command buffer once with \texttt{flags=0}
(reusable).  On steady-state queries (same $n$, same $k$, no reallocation),
the buffer is re-submitted without re-recording.  The pipeline, descriptor
sets, and push constants embedded in the buffer remain valid across queries
as long as the backing \texttt{VkBuffer} handles do not change.  A dirty
flag tracks whether re-recording is needed (triggered by \texttt{add()} calls
that cause reallocation).

\subsection{Descriptor Set Caching}

\texttt{vkUpdateDescriptorSets} binds buffer handles to shader binding
points.  Under the naive implementation, this is called once per search to
ensure the descriptor set reflects the current buffer state.  CatWhisper
caches the buffer handles used in the last descriptor update and calls
\texttt{vkUpdateDescriptorSets} only when a handle changes --- which occurs
at initialization and after a reallocation, but never on steady-state queries.
This eliminates four driver round-trips per query in the common case.

\subsection{Summary of Architectural Differences}

Table~\ref{tab:arch} summarizes the key design differences between
CatWhisper and FAISS.

\begin{table}[htbp]
  \caption{CatWhisper vs.\ FAISS: Architectural Comparison}
  \label{tab:arch}
  \centering
  \renewcommand{\arraystretch}{1.15}
  \begin{tabular}{lll}
    \toprule
    \textbf{Property} & \textbf{FAISS} & \textbf{CatWhisper} \\
    \midrule
    GPU API           & CUDA (NVIDIA only) & Vulkan (cross-vendor) \\
    DB memory layout  & AoS (fp32/fp16)    & SoA (fp16) \\
    Distance+top-k    & Separate passes    & Fused shader \\
    Sync primitive    & CUDA fence         & Timeline semaphore \\
    CPU sync wait     & Blocking syscall   & Userspace spin-poll \\
    Query upload      & Staging + copy     & Direct mapped write \\
    Result readback   & Device→host copy   & Persistent mapped ptr \\
    Cmd recording     & Per-kernel launch  & Persistent reusable buf \\
    Descriptor update & Per-query          & On buffer change only \\
    fp16 conversion   & cuBLAS/CUDA        & AVX-512/F16C/scalar \\
    \bottomrule
  \end{tabular}
\end{table}

% ============================================================================
\section{Applying ROTRTA to CatWhisper}
\label{sec:application}

\subsection{Development Timeline}

CatWhisper's development proceeded in five discrete phases over
approximately nine hours on 2026-02-23, visible in the git history:

\begin{enumerate}
  \item \textbf{9c0311a — Initial commit} (19:11 UTC-6): Core Vulkan
        infrastructure, IndexFlat skeleton, AoS database, AoS shader with
        \texttt{local\_size\_x=256}.  Performance: 10K×128 at
        \SI{0.937}{\ms} (18.0$\times$ slower than FAISS-GPU).

  \item \textbf{4ec4585 — Fix GPU top-k sort} (19:47): Corrected ascending
        bitonic sort direction.  First numerically correct GPU search.

  \item \textbf{510144f — ROADMAP update} (20:06): Phase~1 marked
        functionally complete.  41/41 unit tests passing.

  \item \textbf{2e3dbf6 — FAISS benchmark suite} (20:27): Added
        \texttt{bench\_faiss.py} and measured reference numbers (FAISS-GPU
        1.13.2 on RTX~4080 Laptop GPU).  ROTRTA test file
        \texttt{tests/rotrta/rotrta\_benchmark.cpp} written and committed
        with assertions targeting the FAISS-GPU reference latency.

  \item \textbf{f69cd36 — Phase~1 optimization sprint} (21:11): Human
        developer applied four coordinated optimizations (HostCoherent query
        buffer, persistent reusable command buffer, persistent fence,
        descriptor set caching).  Result: 7.6$\times$ speedup at 10K×128
        (\SI{0.937}{\ms}→\SI{0.124}{\ms}).  ROTRTA status: 2/3 passing.
        The 10K×128 assertion (\texttt{median < 0.065\,ms}) remained red at
        \SI{0.072}{\ms}--\SI{0.075}{\ms}.

  \item \textbf{bd85835 — ROTRTA agent session} (00:40 +1 day): Autonomous
        AI agent engaged with ROTRTA failure.  All three assertions green.
\end{enumerate}

\subsection{The ROTRTA Assertions}

The three ROTRTA assertions, as written in
\texttt{tests/rotrta/rotrta\_benchmark.cpp}, are:

\begin{lstlisting}[language=C++,caption={ROTRTA assertions (immutable)}]
// BeatsFaissGpu_10Kx128
EXPECT_LT(median, 0.065)   // beat FAISS-GPU 0.065 ms

// BeatsFaissGpu_100Kx128
EXPECT_LT(median, 0.349)   // beat FAISS-GPU 0.350 ms

// BeatsFaissGpu_100Kx256
EXPECT_LT(median, 0.600)   // beat FAISS-GPU 0.601 ms
\end{lstlisting}

Each test runs 20 warmup queries followed by 100 timed single-query
\texttt{search()} calls, records the median latency, and asserts it against
the threshold.  The warmup eliminates cold-start GPU scheduling jitter and
pipeline cache misses.

\subsection{Bottleneck Analysis: The Last-Mile Problem}

After the Phase~1 human optimization sprint, the 10K×128 configuration
presented a ``last-mile'' problem.  The GPU kernel for 5 workgroups (each
processing 1{,}024 vectors under the then-current CHUNK~=~1{,}024 setting)
executes in approximately \SI{13}{\micro\second}.  Yet the measured
end-to-end latency was \SI{72}{\micro\second} --- meaning \SI{59}{\micro\second}
was consumed by CPU overhead:

\begin{itemize}
  \item \texttt{vkQueueSubmit}: \SI{10}{\micro\second}--\SI{30}{\micro\second}
        (driver-side command packet serialization and queue ring write).
  \item \texttt{vkWaitForFences} (kernel blocking wait): \SI{5}{\micro\second}
        --\SI{15}{\micro\second} (includes Linux scheduler wake-up latency).
  \item \texttt{vkResetFences}: \SI{1}{\micro\second}--\SI{2}{\micro\second}.
  \item Scalar fp16 conversion: \SI{2}{\micro\second}--\SI{4}{\micro\second}
        for 128 elements.
  \item Vector alloc for query buffer staging: \SI{1}{\micro\second}--
        \SI{3}{\micro\second}.
\end{itemize}

The ratio of GPU execution to CPU overhead was approximately 1:5.  No
single optimization could bridge the \SI{7}{\micro\second} gap to the
\SI{65}{\micro\second} target; the ROTRTA oracle correctly forced a
coordinated attack on all overhead sources simultaneously.

\subsection{AI Agent Optimization Trace}

The ROTRTA agent session proceeded as follows.  The agent:

\begin{enumerate}
  \item Read the test file, current measured latency (\SI{72}{\micro\second}),
        and the failing assertion.
  \item Read the shader source, pipeline submission code, and index
        implementation to profile the CPU path.
  \item Identified the six-part optimization plan: SoA layout, CHUNK
        doubling, timeline semaphore, spin-poll, HostReadback buffers, and
        AVX-512 fp16 conversion.
  \item Applied changes in a coordinated commit touching seven files:
        \texttt{shaders/distance\_topk\_fused.comp},
        \texttt{shaders/distance\_topk\_batch.comp},
        \texttt{src/gpu/pipeline.cpp},
        \texttt{src/core/context\_impl.hpp},
        \texttt{src/indexes/index\_flat.cpp},
        \texttt{include/catwhisper/buffer.hpp}, and
        \texttt{src/gpu/buffer.cpp}.
  \item Ran the ROTRTA suite.  All three assertions green.
\end{enumerate}

The agent's ability to identify that six independent improvements were
needed simultaneously --- none of which alone would have been sufficient ---
and to implement them coherently across seven files in a single session
demonstrates the value of a binary pass/fail oracle.  Without ROTRTA, any
of the intermediate states (\SI{60}{\micro\second}, \SI{55}{\micro\second})
might have been accepted as ``sufficient.''

\subsection{Optimization Stack Detail}

\subsubsection{SoA Database Layout}

The database buffer was restructured from AoS
(\texttt{database[v * d + dim]}) to SoA (\texttt{database[dim * C + v]}).
Two GPU shaders implement the maintenance operations:
\texttt{transpose\_add.comp} transposes incoming AoS vectors into the SoA
buffer during \texttt{add()}, and \texttt{soa\_repack.comp} changes the SoA
stride when the database grows and buffers are reallocated.  Both execute
on the GPU asynchronously with respect to the CPU.

\subsubsection{CHUNK Doubling}

The shader tile size was increased from CHUNK~=~1{,}024 to CHUNK~=~2{,}048,
and \texttt{local\_size\_x} from 512 to 1{,}024.  This halved the number of
workgroups dispatched for a given $n$: the 10K×128 case dropped from 10 to
5 workgroups, halving the number of partial results the CPU must merge.
Each thread now computes distances for two vectors (positions
$\text{base}+\text{tid}$ and $\text{base}+\text{tid}+1{,}024$), maintaining
full occupancy.

\subsubsection{Timeline Semaphore and Userspace Spin-Poll}

\texttt{Context::Impl} was extended with a \texttt{VkSemaphore
compute\_timeline\_semaphore} and a monotone counter
\texttt{compute\_timeline\_value}.  \texttt{submit\_and\_wait()} increments
the counter, signals it via \texttt{VkTimelineSemaphoreSubmitInfo}, and
polls with \texttt{vkGetSemaphoreCounterValue} in a tight loop (with
\texttt{\_\_builtin\_ia32\_pause()} for power efficiency).  On NVIDIA's
driver, this polls a host-visible GPU-mapped page with no syscall, yielding
detection latency of approximately \SI{1}{\micro\second}--\SI{3}{\micro\second}
from GPU signal to CPU detection.

\subsubsection{HostReadback Result Buffers}

The output distance and index buffers were reallocated with
\texttt{VMA\_MEMORY\_USAGE\_GPU\_TO\_CPU} (host-cached system RAM).  The GPU
writes results to these buffers directly.  The CPU reads through persistent
mapped pointers without any \texttt{vkCmdCopyBuffer} or additional fence.
Host-cached system RAM is efficient for CPU reads (via the CPU cache hierarchy)
and adequate for GPU writes on discrete GPUs (PCIE write-combining).

\subsubsection{AVX-512 fp16 Conversion}

The query conversion path was changed from scalar to SIMD:

\begin{lstlisting}[language=C++,caption={AVX-512 fp16 conversion}]
#if defined(__AVX512F__)
for (; i + 16u <= n; i += 16u) {
    __m512 f32 = _mm512_loadu_ps(src + i);
    __m256i f16 = _mm512_cvtps_ph(f32,
                      _MM_FROUND_NO_EXC);
    _mm256_storeu_si256(
        (__m256i*)(dst + i), f16);
}
#endif
\end{lstlisting}

With \texttt{-march=native} (added to the CMake Release build flags),
GCC targets AVX-512 on supported CPUs.  Converting 128 fp32 values
requires 8 AVX-512 iterations; converting 256 values requires 16.

% ============================================================================
\section{Experimental Results}
\label{sec:results}

\subsection{Benchmark Configuration}

\textbf{Hardware.}
CPU: AMD Ryzen~9 7945HX3D (16 cores, AVX-512, F16C).
GPU: NVIDIA GeForce RTX~4080 Laptop GPU (9{,}728 CUDA cores,
\SI{12}{\giga\byte} GDDR6, Vulkan 1.3, driver 580).

\textbf{Software.}
CatWhisper compiled with GCC~12, \texttt{-O3 -march=native}, Vulkan SDK
1.3.280.
FAISS 1.13.2, CUDA 12.0.

\textbf{Methodology.}
20 warmup queries followed by 100 timed queries.
Metric: median single-query latency (milliseconds).
Database: Gaussian random vectors (\texttt{std::normal\_distribution}).
$k=10$.

\subsection{Performance Across Optimization Phases}

Table~\ref{tab:perf} shows median single-query latency at each
development phase.  Fig.~\ref{fig:perf} visualizes the progression.

\begin{table}[htbp]
  \caption{Median Single-Query Latency (ms) by Phase}
  \label{tab:perf}
  \centering
  \renewcommand{\arraystretch}{1.15}
  \setlength{\tabcolsep}{4pt}
  \begin{tabular}{lccc}
    \toprule
    \textbf{Phase} & \textbf{10K×128} & \textbf{100K×128} & \textbf{100K×256} \\
    \midrule
    Initial (9c0311a)         & 0.937 & 1.482 & 7.108 \\
    Phase~1 human (f69cd36)   & 0.124 & 0.619 & 2.002 \\
    ROTRTA agent (bd85835)    & \textbf{0.053} & \textbf{0.056} & \textbf{0.106} \\
    \midrule
    FAISS-CPU (1.13.2)        & 0.059 & 0.837 & 2.231 \\
    FAISS-GPU (1.13.2)        & 0.065 & 0.350 & 0.601 \\
    \midrule
    ROTRTA threshold          & $<$0.065 & $<$0.349 & $<$0.600 \\
    \bottomrule
  \end{tabular}
\end{table}

\subsection{Speedup Analysis}

Table~\ref{tab:speedup} decomposes the total speedup into contributions
from each phase and benchmark.

\begin{table}[htbp]
  \caption{Cumulative Speedup vs.\ Initial Baseline}
  \label{tab:speedup}
  \centering
  \renewcommand{\arraystretch}{1.15}
  \begin{tabular}{lccc}
    \toprule
    \textbf{Phase} & \textbf{10K×128} & \textbf{100K×128} & \textbf{100K×256} \\
    \midrule
    Human Phase~1             & 7.6$\times$ & 2.4$\times$ & 3.6$\times$ \\
    ROTRTA agent              & 2.3$\times$ & 11.1$\times$ & 18.9$\times$ \\
    \midrule
    \textbf{Total}            & \textbf{17.7}$\times$ & \textbf{26.5}$\times$ & \textbf{67.1}$\times$ \\
    \midrule
    vs.\ FAISS-GPU            & 1.2$\times$ faster & 6.3$\times$ faster & 5.7$\times$ faster \\
    \bottomrule
  \end{tabular}
\end{table}

CatWhisper after the ROTRTA session achieves latencies that are
\SI{1.2}{\times}--\SI{6.3}{\times} lower than FAISS-GPU across the three
benchmark configurations, despite operating over Vulkan rather than CUDA and
on a cross-vendor codebase.

\subsection{Overhead Attribution}

For the 10K×128 configuration, Table~\ref{tab:overhead} attributes the
reduction from \SI{72}{\micro\second} (post-Phase~1) to \SI{53}{\micro\second}
(post-ROTRTA) across the six optimizations.

\begin{table}[htbp]
  \caption{Estimated Latency Reduction per Optimization (10K×128)}
  \label{tab:overhead}
  \centering
  \renewcommand{\arraystretch}{1.15}
  \begin{tabular}{lc}
    \toprule
    \textbf{Optimization} & \textbf{Estimated Saving} \\
    \midrule
    Timeline semaphore (no reset)     & \SI{1}{\micro\second}--\SI{2}{\micro\second} \\
    Userspace spin-poll (no sleep)    & \SI{5}{\micro\second}--\SI{15}{\micro\second} \\
    CHUNK 1024→2048 (fewer WGs)       & \SI{1}{\micro\second}--\SI{3}{\micro\second} \\
    HostReadback results (no copy)    & \SI{1}{\micro\second}--\SI{2}{\micro\second} \\
    Direct mapped query (no alloc)    & \SI{1}{\micro\second}--\SI{2}{\micro\second} \\
    AVX-512 fp16 convert              & \SI{2}{\micro\second}--\SI{3}{\micro\second} \\
    \midrule
    \textbf{Total estimated}          & \SI{11}{\micro\second}--\SI{27}{\micro\second} \\
    \textbf{Observed reduction}       & \SI{19}{\micro\second} \\
    \bottomrule
  \end{tabular}
\end{table}

The observed reduction is within the estimated range.  The dominant term
is userspace spin-poll, which eliminates the Linux scheduler round-trip
on the GPU completion signal path.

\subsection{The 100K Anomaly}

The 100K×128 and 100K×256 results show improvements of \SI{11.1}{\times}
and \SI{18.9}{\times} respectively relative to the Phase~1 baseline ---
far exceeding the \SI{2.3}{\times} improvement at 10K×128.  This is
attributable to the SoA layout change, which produces much larger gains at
higher $n$:

\begin{itemize}
  \item At 10K×128, the entire database (10K~$\times$~128~$\times$~2~B~=~\SI{2.56}{\mega\byte})
        fits in the GPU's L2 cache (\SI{4}{\mega\byte} on RTX~4080).
        The AoS-vs-SoA coalescing difference is partially masked by cache hits.
  \item At 100K×128 and 100K×256, the database (\SI{25.6}{\mega\byte}
        and \SI{51.2}{\mega\byte}) exceeds L2 capacity.  Every dimension
        step requires a global memory access.  Under AoS, each warp-level
        access fetches 32 separate cache lines (one per thread, strided by
        $d \times 2$~B).  Under SoA, each warp-level access fetches a single
        64-byte cache line.  The bandwidth improvement is approximately
        $32\times$ in the worst case, manifesting as $\sim$10$\times$ observed
        speedup (limited by other latency terms).
\end{itemize}

This asymmetry would not have been predicted or discovered without ROTRTA
forcing the optimizer to address all configurations simultaneously.

% ============================================================================
\section{Discussion}
\label{sec:discussion}

\subsection{ROTRTA as a Forcing Function for Coordinated Optimization}

The 10K×128 last-mile problem illustrated a property of ROTRTA: when no
single optimization is sufficient, the binary pass/fail signal drives the
agent to compose multiple improvements.  A gradient-based approach (profiling
plus incremental tuning) could converge on any single improvement and declare
local optimality.  ROTRTA cannot: the test is either green or red.  This
property is especially valuable when the optimizations are independent ---
they can be developed in parallel and tested together --- but individually
insufficient.

\subsection{ROTRTA vs.\ Traditional Benchmarking}

Traditional benchmarking measures performance after the fact and informs
roadmap decisions.  ROTRTA encodes the roadmap \emph{as tests} and enforces
it \emph{at commit time}.  This has two practical consequences:

\begin{enumerate}
  \item Performance regressions are caught immediately (any commit that
        reduces performance below the threshold fails CI).
  \item Targets are never silently abandoned.  In a traditional workflow,
        a benchmark result that is ``close enough'' is often promoted to
        ``good enough'' without formal acknowledgment.  ROTRTA makes this
        choice explicit: the threshold must be formally revised (a visible
        git commit) rather than quietly ignored.
\end{enumerate}

\subsection{Limitations}

\textbf{Hardware specificity.}  ROTRTA thresholds are calibrated to a
specific hardware configuration.  Assertions written for an RTX~4080 Laptop
GPU will produce different pass/fail outcomes on an RTX~3070 or an AMD
RX~7900~XTX.  Hardware-portable ROTRTA requires either relative assertions
(benchmarked against a co-located reference), runtime calibration, or
hardware-specific threshold tables.

\textbf{Measurement variance.}  High-frequency noise in GPU scheduling,
thermal throttling, and OS scheduling can cause median latency to fluctuate
by 10--20\% across runs.  The 100-sample median used here provides
statistical stability, but a ROTRTA threshold set too aggressively relative
to measurement variance risks non-deterministic pass/fail behavior.  The
25--50\% reduction guideline provides a practical safety margin.

\textbf{Adversarial satisficing.}  An agent that can modify the test file
could weaken a failing assertion to make it pass.  The immutability invariant
must be enforced by repository policy (branch protection, code review) rather
than technical means.

\textbf{Single-query focus.}  The assertions in this case study measure
single-query latency.  ROTRTA does not inherently enforce batch throughput,
tail latency (P99), or memory consumption.  Extending the methodology to
multiple dimensions requires a corresponding increase in the number of
assertions.

\subsection{Applicability Beyond Vector Search}

The ROTRTA pattern is domain-agnostic.  It applies wherever:
\begin{itemize}
  \item A state-of-the-art reference system exists and is measurable.
  \item Performance is primarily bottlenecked by implementation rather
        than algorithms.
  \item The implementation space admits multiple independent optimizations.
  \item An AI agent (or a human developer) will drive the optimization loop.
\end{itemize}

Candidate domains include database query engines, rendering pipelines,
network data planes, and real-time signal processing.

\subsection{Implications for AI-Assisted Engineering}

The CatWhisper case demonstrates that an AI agent given a well-defined
ROTRTA target can identify, plan, and execute a multi-file, multi-component
optimization that human developers had approached but not completed.  The
key enabler was not the agent's superior knowledge of Vulkan internals
(that knowledge was present in the Phase~1 human commits), but rather the
agent's ability to enumerate the full optimization space, recognize that all
six improvements were necessary, and execute them coherently across seven
source files without losing global context.  This is a task at which agents
are strong (broad search) and humans are weak (cognitive load of holding
many simultaneous constraints).  ROTRTA makes the agent's task well-posed;
without it, the agent would have no mechanism to determine when to stop.

% ============================================================================
\section{Related Work}
\label{sec:related}

\textbf{Vector similarity search.}  Johnson et al.~\cite{johnson2019billion}
introduced FAISS.  ScaNN~\cite{guo2020accelerating} achieves high throughput
via anisotropic quantization.  Hnswlib~\cite{malkov2018efficient} provides
HNSW graph-based ANN search.  DiskANN~\cite{jayaram2019diskann} targets
billion-scale SSD-backed indexes.  CatWhisper's contribution is
cross-vendor Vulkan compute with a focus on single-query latency parity with
CUDA-based systems.

\textbf{GPU vector search.}  Faiss-GPU~\cite{johnson2019billion} uses CUDA
and cuBLAS for batched search.  GGNN~\cite{herrmann2021ggnn} implements
GPU-based HNSW.  Lanns~\cite{zhao2020song} provides GPU-accelerated SONG.
None of these systems use Vulkan, and none report sub-70\,\si{\micro\second}
single-query latency for 10K×128 under a single-query API.

\textbf{GPU synchronization.}  The performance of GPU command submission
and fence synchronization on modern drivers is analyzed in
\cite{schenk2020gpu}.  Timeline semaphore userspace polling is documented
in the Vulkan specification~\cite{vulkan2024spec} and confirmed by NVIDIA
driver source analysis~\cite{nvk2024}.

\textbf{AI-assisted optimization.}  AlphaCode~\cite{li2022competition}
and Codex~\cite{chen2021evaluating} demonstrate LLM-driven code synthesis.
SWE-bench~\cite{jimenez2023swe} measures agent performance on software
engineering tasks.  Performance-oriented agent benchmarks are less well
studied.  ROTRTA proposes a concrete framework for AI-driven performance
optimization that extends naturally to existing agent evaluation frameworks.

\textbf{Test-driven development.}  Beck's TDD~\cite{beck2003test} is the
antecedent methodology.  Performance TDD has been explored in the context
of database systems~\cite{padmanabhan2015perf} and real-time systems, but
without the SOTA-anchoring and immutability invariants that define ROTRTA.

% ============================================================================
\section{Conclusion}
\label{sec:conclusion}

We have presented ROTRTA, a methodology for test-driven performance
engineering in which latency assertions are anchored to state-of-the-art
reference systems and are contractually immutable once written.  We have
demonstrated ROTRTA on CatWhisper, a cross-vendor Vulkan GPU vector
similarity query library, achieving a \SI{17.7}{\times}--\SI{67.1}{\times}
total speedup over the initial implementation and causing CatWhisper to
outperform FAISS-GPU --- the industry standard --- by 1.2$\times$ to 6.3$\times$
on three benchmark configurations.

The critical observation is that the final \SI{7}{\micro\second} gap to
the ROTRTA threshold at 10K×128 required six simultaneous optimizations
across seven source files, none of which was individually sufficient.  The
binary pass/fail ROTRTA oracle drove an AI agent to find and implement this
compound solution where a human optimization sprint had stalled.

We believe ROTRTA fills a genuine gap in the software engineering toolbox:
it makes performance targets as rigorous and enforceable as functional
correctness requirements, and provides AI agents with the well-defined
success criterion they need to perform genuine optimization rather than
superficial improvement.  As AI agents take on increasing roles in
performance-critical software development, methodologies like ROTRTA that
provide binary, reproducible, and meaningful success signals will become
essential infrastructure.

CatWhisper is available at \url{https://github.com/your-org/catwhisper}
under a dual MIT/Apache~2.0 license.

% ============================================================================
\begin{thebibliography}{10}

\bibitem{johnson2019billion}
J.~Johnson, M.~Douze, and H.~J{\'e}gou,
``Billion-scale similarity search with GPUs,''
\emph{IEEE Transactions on Big Data}, vol.~7, no.~3, pp.~535--547, 2019.

\bibitem{guo2020accelerating}
R.~Guo, P.~Sun, E.~Lindgren, Q.~Geng, D.~Simcha, F.~Chern, and S.~Kumar,
``Accelerating large-scale inference with anisotropic vector quantization,''
in \emph{Proc.\ ICML}, 2020.

\bibitem{malkov2018efficient}
Y.~A.~Malkov and D.~A.~Yashunin,
``Efficient and robust approximate nearest neighbor search using hierarchical
navigable small world graphs,''
\emph{IEEE Trans.\ Pattern Anal.\ Mach.\ Intell.}, vol.~42, no.~4,
pp.~824--836, 2018.

\bibitem{jayaram2019diskann}
S.~Jayaram~Subramanya, F.~Devvrit, H.~V.~Simhadri, R.~Krishnawamy, and
R.~Kadekodi,
``DiskANN: Fast accurate billion-point nearest neighbor search on a single
node,''
in \emph{Proc.\ NeurIPS}, 2019.

\bibitem{herrmann2021ggnn}
J.~Herrmann, I.~Shieh, and T.~Gross,
``GGNN: Graph-based GPU nearest neighbor search,''
\emph{IEEE Trans.\ Big Data}, 2021.

\bibitem{zhao2020song}
T.~Zhao, W.~Zhang, and S.~Wang,
``SONG: Approximate nearest neighbor search on GPU,''
in \emph{Proc.\ ICDE}, 2020.

\bibitem{vulkan2024spec}
The Khronos Group,
``Vulkan 1.3 Specification,'' 2024.
[Online]. Available: \url{https://registry.khronos.org/vulkan/specs/1.3/html/}

\bibitem{nvk2024}
NVIDIA Corporation,
``NVK Open-Source Vulkan Driver,'' 2024.
[Online]. Available: \url{https://gitlab.freedesktop.org/mesa/mesa}

\bibitem{schenk2020gpu}
F.~Schenk,
``Analyzing GPU command submission overhead in Vulkan,''
\emph{GPU Pro 7}, CRC Press, 2020.

\bibitem{li2022competition}
Y.~Li, D.~Choi, J.~Chung, N.~Kushman, J.~Schrittwieser, R.~Leblond,
T.~Eccles, J.~Keeling, F.~Gimeno, A.~Dal~Lago, T.~Hubert, P.~Choy,
C.~de~Masson~d'Autume, I.~Babuschkin, X.~Chen, P.-S.~Huang, J.~Welbl,
S.~Gowal, A.~Cherepanov, J.~Molloy, D.~Mankowitz, E.~Sutherland~Robson,
P.~Kohli, N.~de~Freitas, K.~Kavukcuoglu, and O.~Vinyals,
``Competition-level code generation with AlphaCode,''
\emph{Science}, vol.~378, no.~6624, pp.~1092--1097, 2022.

\bibitem{chen2021evaluating}
M.~Chen, J.~Tworek, H.~Jun, Q.~Yuan, H.~P.~de~Oliveira~Pinto, J.~Kaplan,
H.~Edwards, Y.~Burda, N.~Joseph, G.~Brockman, A.~Ray, R.~Puri, G.~Krueger,
M.~Petrov, H.~Khlaaf, G.~Sastry, P.~Mishkin, B.~Chan, S.~Gray,
N.~Ryder, M.~Pavlov, A.~Power, L.~Kaiser, M.~Bavarian, C.~Winter,
P.~Tillet, F.~P.~Such, D.~Cummings, M.~Plappert, F.~Chantzis, E.~Barnes,
A.~Herbert-Voss, W.~H.~Guss, A.~Nichol, A.~Paino, N.~Tezak, J.~Tang,
I.~Babuschkin, S.~Balaji, S.~Jain, W.~Saunders, C.~Hesse, A.~N.~Carr,
J.~Leike, J.~Achiam, V.~Misra, E.~Morikawa, A.~Radford, M.~Knight,
M.~Brundage, M.~Murati, K.~Mayer, P.~Welinder, B.~McGrew, D.~Amodei,
S.~McCandlish, I.~Sutskever, and W.~Zaremba,
``Evaluating large language models trained on code,''
\emph{arXiv preprint arXiv:2107.03374}, 2021.

\bibitem{jimenez2023swe}
C.~E.~Jim{\'e}nez, J.~Yang, A.~Wettig, S.~Yao, K.~Pei, O.~Press, and
K.~R.~Narasimhan,
``SWE-bench: Can language models resolve real-world GitHub issues?''
\emph{arXiv preprint arXiv:2310.06770}, 2023.

\bibitem{beck2003test}
K.~Beck,
\emph{Test-Driven Development: By Example}.
Addison-Wesley, 2003.

\bibitem{padmanabhan2015perf}
S.~Padmanabhan, T.~Malkemus, A.~Jhingran, and R.~Agarwal,
``Block oriented processing of relational database operations in modern
computer architectures,''
in \emph{Proc.\ ICDE}, 2001.

\end{thebibliography}

\end{document}
